Mzdová stratifikace podle pohlaví a~vzdělání ukazuje, že \ac{GPG} v~ČR není homogenní: je výrazně větší u~středoškolsky vzdělaných pracovníků než u~vysokoškolsky vzdělaných, a~ve zpracovatelském průmyslu výrazněji než v~ICT nebo financích.
Tento vzorec odpovídá segmentaci trhu práce: segmenty, kde jsou ženy koncentrovány (péče, administrativa, maloobchod), mají zároveň slabší odborovou organizaci a~nižší podíl zaměstnanců pokrytých \ac{KS} --- dvojí nevýhoda pohlaví a~institucionální dezorganizace se tak vzájemně zesilují.
Výzkum Eurofound (2020) dokumentuje, že v~odvětvích, kde jsou odvětvové \ac{KS} s~transparentní tarifní soustavou, je mezera mezi mzdami žen a~mužů statisticky signifikantně nižší --- efekt je nezávislý na odvětvové a~profesní struktuře.
Systémové řešení \ac{GPG} v~ČR proto předpokládá nejen legislativní transparentnost odměňování, ale především posílení kolektivního pokrytí v~feminizovaných odvětvích, kde je vyjednávací síla žen systematicky nejnižší.

Mzdová stratifikace podle pohlaví a~vzdělání ukazuje, že \ac{GPG} v~ČR není homogenní: je výrazně větší u~středoškolsky vzdělaných pracovníků než u~vysokoškolsky vzdělaných, a~ve zpracovatelském průmyslu výrazněji než v~ICT nebo financích.
Tento vzorec odpovídá segmentaci trhu práce: segmenty, kde jsou ženy koncentrovány (péče, administrativa, maloobchod), mají zároveň slabší odborovou organizaci a~nižší podíl zaměstnanců pokrytých \ac{KS} --- dvojí nevýhoda pohlaví a~institucionální dezorganizace se tak vzájemně zesilují.
Výzkum Eurofound (2020) dokumentuje, že v~odvětvích, kde jsou odvětvové \ac{KS} s~transparentní tarifní soustavou, je mezera mezi mzdami žen a~mužů statisticky signifikantně nižší --- efekt je nezávislý na odvětvové a~profesní struktuře.
Systémové řešení \ac{GPG} v~ČR proto předpokládá nejen legislativní transparentnost odměňování, ale především posílení kolektivního pokrytí v~feminizovaných odvětvích, kde je vyjednávací síla žen systematicky nejnižší.

Mzdová stratifikace podle pohlaví a~vzdělání ukazuje, že \ac{GPG} v~ČR není homogenní: je výrazně větší u~středoškolsky vzdělaných pracovníků než u~vysokoškolsky vzdělaných, a~ve zpracovatelském průmyslu výrazněji než v~ICT nebo financích.
Tento vzorec odpovídá segmentaci trhu práce: segmenty, kde jsou ženy koncentrovány (péče, administrativa, maloobchod), mají zároveň slabší odborovou organizaci a~nižší podíl zaměstnanců pokrytých \ac{KS} --- dvojí nevýhoda pohlaví a~institucionální dezorganizace se tak vzájemně zesilují.
Výzkum Eurofound (2020) dokumentuje, že v~odvětvích, kde jsou odvětvové \ac{KS} s~transparentní tarifní soustavou, je mezera mezi mzdami žen a~mužů statisticky signifikantně nižší --- efekt je nezávislý na odvětvové a~profesní struktuře.
Systémové řešení \ac{GPG} v~ČR proto předpokládá nejen legislativní transparentnost odměňování, ale především posílení kolektivního pokrytí v~feminizovaných odvětvích, kde je vyjednávací síla žen systematicky nejnižší.

Mzdová stratifikace podle pohlaví a~vzdělání ukazuje, že \ac{GPG} v~ČR není homogenní: je výrazně větší u~středoškolsky vzdělaných pracovníků než u~vysokoškolsky vzdělaných, a~ve zpracovatelském průmyslu výrazněji než v~ICT nebo financích.
Tento vzorec odpovídá segmentaci trhu práce: segmenty, kde jsou ženy koncentrovány (péče, administrativa, maloobchod), mají zároveň slabší odborovou organizaci a~nižší podíl zaměstnanců pokrytých \ac{KS} --- dvojí nevýhoda pohlaví a~institucionální dezorganizace se tak vzájemně zesilují.
Výzkum Eurofound (2020) dokumentuje, že v~odvětvích, kde jsou odvětvové \ac{KS} s~transparentní tarifní soustavou, je mezera mezi mzdami žen a~mužů statisticky signifikantně nižší --- efekt je nezávislý na odvětvové a~profesní struktuře.
Systémové řešení \ac{GPG} v~ČR proto předpokládá nejen legislativní transparentnost odměňování, ale především posílení kolektivního pokrytí v~feminizovaných odvětvích, kde je vyjednávací síla žen systematicky nejnižší.

\input{texparts/python/gender_wage_stratification}



